\chapter{FUNDAMENTAÇÃO TEÓRICA E TRABALHOS RELACIONADOS}
\label{ch:fundamentacao}

Este capítulo apresenta os conceitos e trabalhos que fundamentam a análise de transformação de notícias em cadeias de reescrita. A discussão parte das relações entre deriva informacional, desinformação e enquadramento, situando os limites da classificação de veracidade para descrever mudanças entre versões. Em seguida, examina o condicionamento de agentes por perspectivas e as representações semânticas e afetivas usadas na comparação de textos. Por fim, discute a avaliação automática e sintetiza o posicionamento do estudo em relação aos trabalhos mais próximos.

\section{TRANSFORMAÇÃO TEXTUAL, DESINFORMAÇÃO E ENQUADRAMENTO}
\label{sec:evolucao-enquadramento}

Notícias falsas, rumores e desinformação podem compartilhar conteúdos e canais de circulação, mas diferem quanto à forma, ao estado de verificação e, em algumas definições, à intenção de quem produz ou divulga a mensagem. \citeonline{lazer2018science} situam as notícias falsas como informações fabricadas que imitam a forma do conteúdo jornalístico, sem adotar seus processos editoriais e normas de produção. Rumores, por sua vez, não constituem uma classe necessariamente falsa: uma alegação pode circular enquanto sua veracidade permanece incerta e ser posteriormente confirmada ou refutada. A revisão de \citeonline{yang2025rethink} organiza a detecção de rumores na era dos modelos de linguagem em níveis de cognição, interação e comportamento, abrangendo o que participantes expressam, suas relações e suas ações durante a propagação.

Esses objetos exigem distinguir unidades de análise. A veracidade factual diz respeito à correspondência entre uma afirmação e evidências sobre o mundo externo. A crença corresponde à aceitação dessa afirmação, a exposição registra o contato com o conteúdo e o compartilhamento constitui uma ação de circulação. Os padrões de difusão observados por \citeonline{vosoughi2018spread}, por exemplo, caracterizam a propagação de notícias verdadeiras e falsas, mas não descrevem como uma mensagem muda durante esse percurso. Para examinar a transformação textual, é necessário comparar o conteúdo das versões transmitidas.

Neste TCC, \emph{deriva informacional} designa a mudança mensurável entre versões de um conteúdo, incluindo alterações de tema, subtópicos, entidades, relações e enquadramento. Essa comparação não determina qual versão é verdadeira. Uma reescrita pode afastar-se da original e permanecer factualmente correta, assim como uma alteração factual decisiva pode ocorrer em textos semanticamente próximos. Também se distingue deriva de \emph{contradição interna}, entendida como incompatibilidade entre proposições de uma mesma versão. Dois textos podem divergir entre si e, ainda assim, ser internamente coerentes. A \emph{deriva afetiva}, por sua vez, descreve a mudança no perfil emocional estimado do texto, sem representar uma mudança comprovada no estado psicológico de seus leitores ou autores.

A possibilidade de transformação durante a transmissão antecede o uso de grandes modelos de linguagem, ou \textit{Large Language Models} (LLMs). \citeonline{zhang2013rumor} incorporam a modificação da mensagem à modelagem da evolução de rumores, considerando indivíduos que podem encaminhar ou alterar o conteúdo. Em uma abordagem empírica, \citeonline{guo2021truth} apresentam o conjunto Fake News Evolution, organizado em estágios de notícia verdadeira, notícia falsa e notícia falsa evoluída. Os autores examinam propriedades como similaridade textual, palavras-chave e sentimento, aproximando evolução de conteúdo e detecção. Esses antecedentes tratam a mensagem como um objeto que pode ser modificado ao circular.

O conceito de enquadramento permite caracterizar parte dessas mudanças. Para \citeonline{entman1993framing}, enquadrar envolve selecionar aspectos de uma realidade percebida e torná-los salientes, favorecendo determinadas definições de problema, interpretações causais, avaliações ou recomendações. Assim, uma reescrita pode conservar o acontecimento principal e seus participantes enquanto altera a responsabilização, a justificativa de uma ação ou a importância atribuída às suas consequências. A proximidade temática não garante a preservação dessas relações. Inversamente, diferenças de vocabulário podem corresponder a uma paráfrase que mantém o sentido.

A geração iterativa por LLMs permite investigar esse processo em condições controladas. \citeonline{mohamed2025broken} utilizam cadeias de tradução para examinar o que ocorre quando a saída de uma etapa alimenta a seguinte. Os autores observam acúmulo de distorção associado ao idioma e à complexidade da cadeia e estudam estratégias de \textit{prompt} para mitigá-lo. O trabalho motiva a observação de versões intermediárias, além da comparação final, embora sua tarefa de tradução seja distinta da reescrita condicionada por perspectivas adotada neste TCC.

Uma aproximação mais direta com a transformação de notícias aparece em \citeonline{liu2025stepwise}. Os autores propõem o FUSE, uma simulação em rede com agentes que propagam, comentam, verificam ou observam mensagens durante a evolução de notícias verdadeiras para falsas. O FUSE-EVAL avalia o desvio em dimensões linguísticas e semânticas e é comparado a julgamentos humanos. O interesse pela transformação gradual e pela avaliação multidimensional aproxima esse trabalho do presente estudo. A diferença está no alvo e no desenho: aqui, o foco é o afastamento de uma referência textual ao longo de sequências controladas de perspectivas.

A classificação automática de veracidade responde a uma questão diferente dessa comparação. \citeonline{su2024adapting} avaliam detectores em combinações de notícias escritas por humanos, geradas por máquinas e parafraseadas por máquinas, mostrando assimetrias de generalização entre condições. \citeonline{wu2024sheepdog} investigam ataques de estilo produzidos por LLMs contra detectores de notícias falsas. Esses estudos mostram por que o modo de geração e a apresentação do texto precisam ser considerados ao interpretar a classificação. Já \citeonline{whitehouse2022knowledge} examinam a integração de conhecimento externo à detecção, com benefícios dependentes da pertinência e da atualidade da base. Uma classe atribuída ao documento não localiza, por si só, as alterações entre versões, e a comparação entre versões não substitui a consulta a evidências factuais.

Para descrever uma cadeia, duas referências são complementares: a versão imediatamente anterior e o texto original. A primeira localiza a mudança introduzida em uma passagem. A segunda descreve o afastamento alcançado em relação à origem. Uma retransmissão com pouca mudança local pode conservar uma interpretação acrescentada anteriormente. Essa distinção, combinada à análise de conteúdo e enquadramento, orienta o problema de mensuração desenvolvido neste trabalho.

\section{AGENTES LLM E PERSPECTIVAS}
\label{sec:agentes-personas}

Um agente LLM pode ser entendido, neste contexto, como uma configuração computacional que combina um modelo de linguagem, instruções, estado disponível e regras de interação. Uma persona é o condicionamento textual destinado a orientar a saída para um papel ou perspectiva. Seu uso permite investigar como instruções diferentes se relacionam às respostas produzidas, mas a consistência desse condicionamento precisa ser examinada.

\citeonline{chuang2024opinion} empregam populações de agentes LLM para simular dinâmicas de opinião. Em seus experimentos, os agentes apresentam tendências próprias do modelo em direção a informações corretas e ao consenso. A indução de viés de confirmação por \textit{prompts} produz padrões de fragmentação. Esse resultado indica que a dinâmica observada depende tanto da instrução quanto de predisposições do modelo. A perspectiva solicitada é, portanto, uma das condições que participam da geração da resposta.

A estabilidade dos perfis é examinada por \citeonline{frisch2024interaction}, que encontram diferentes graus de consistência de personalidade e de alinhamento linguístico em tarefas individuais e interativas. Em cadeias de reescrita, essa preocupação torna relevante observar a relação entre a instrução atual e as características da mensagem recebida. A presença de uma persona no \textit{prompt} não basta para estabelecer que todas as saídas expressem o papel de maneira uniforme.

Outras simulações ajudam a delimitar o objeto de análise. \citeonline{liu2024skepticism} investigam mudanças de atitude diante de notícias falsas com agentes dotados de personalidade, memória e reflexão. \citeonline{wang2025echo} analisam polarização em redes simuladas, enquanto \citeonline{liu2025mosaic} modelam disseminação e regulação de conteúdo por ações como publicar, compartilhar e moderar. Esses trabalhos exploram atitudes e comportamentos atribuídos a agentes. No presente estudo, a unidade de interesse é a versão textual produzida em cada passagem de uma cadeia.

A avaliação de agentes com papéis também exige critérios próprios. \citeonline{chen2025rpa} organizam dimensões e métodos de avaliação relacionados, entre outros aspectos, à fidelidade e à consistência dos perfis. Essas questões fundamentam a cautela ao interpretar as perspectivas como condições de geração, sem equivalência com personalidades humanas. As saídas de uma cadeia refletem a combinação entre modelo, instrução e texto recebido. A comparação entre sequências deve considerar esse conjunto de condições, bem como as dependências introduzidas pelo desenho experimental.

\section{REPRESENTAÇÃO SEMÂNTICA E AFETIVA}
\label{sec:representacao-afetiva}

A comparação entre versões exige representar aspectos do significado para além da coincidência de palavras. Embeddings de sentenças projetam trechos em vetores cuja proximidade pode ser estimada por medidas como a similaridade de cosseno. \citeonline{reimers2019sbert} propõem o Sentence-BERT (SBERT), que adapta modelos do tipo BERT com estruturas siamesas e de triplas para produzir representações eficientes em tarefas de similaridade e agrupamento. Essa abordagem fundamenta a comparação vetorial de elementos textuais, como expressões que descrevem temas ou relações.

\citeonline{wang2020minilm} apresentam uma estratégia de destilação de autoatenção para comprimir modelos Transformer. O trabalho fundamenta a família MiniLM, da qual deriva o encoder utilizado pelo comparador Cluster do STDI. A representação de sentenças deve ser compreendida em conjunto com a contribuição do SBERT, pois a arquitetura compacta, por si só, não descreve o ajuste específico para similaridade. A identificação do encoder e as regras de comparação adotadas no estudo são apresentadas no Capítulo \ref{ch:metodos}.

Representações vetoriais permitem reconhecer proximidade entre expressões distintas, mas a similaridade global pode ocultar mudanças localizadas. Duas versões podem tratar do mesmo assunto e divergir quanto a quem realizou uma ação, à sua negação ou a uma quantidade. A decomposição do texto em elementos comparáveis busca tornar essas diferenças examináveis. Sua utilidade depende também da representação intermediária: uma informação preservada no texto pode estar ausente da estrutura extraída. Por isso, a interpretação de uma diferença exige considerar o texto, a extração e o procedimento que os compara.

Além do conteúdo proposicional, reescritas podem modificar o tom afetivo. O modelo de valência, ativação e dominância (VAD) representa emoções em dimensões contínuas. A valência indica quão positivo ou negativo é o conteúdo afetivo, a ativação descreve o grau de excitação ou intensidade e a dominância expressa a percepção de controle ou poder. O EmoBank, apresentado por \citeonline{buechel2017emobank}, oferece anotações dimensionais de sentenças e distingue perspectivas de emoção expressa pelo autor e evocada no leitor. Essa separação é pertinente ao uso de escores automáticos para descrever o tom de uma mensagem.

Há diferentes estratégias para estimar essas dimensões. No nível lexical, \citeonline{mohammad2018vad} apresentam avaliações humanas de VAD para mais de vinte mil palavras em inglês. O NRC VAD v2 amplia o recurso com avaliações de palavras e expressões \cite{mohammad2025vadv2,mohammad2025breaking}. As traduções do NRC v1, incluindo a versão em português, foram produzidas automaticamente a partir dos termos ingleses \cite{mohammad2018vadresource}. Já \citeonline{buechel2020memolon} combinam recursos de uma língua de origem, tradução de palavras e representações vetoriais para gerar léxicos afetivos em 91 idiomas, no recurso MEmoLon. A disponibilidade de um léxico em português não equivale a uma coleta de avaliações humanas de documentos nesse idioma.

Para estimar o perfil de um documento por um léxico, é necessário agregar os valores dos termos reconhecidos. O resultado depende da cobertura e da forma de reconhecer palavras e expressões. Uma média lexical não modela, por si só, o escopo da negação, quem expressa uma avaliação ou a relação entre trechos. A avaliação lexical dos recursos não estabelece automaticamente a validade desse escore agregado por notícia.

Modelos neurais constituem outra estratégia. A ficha oficial de \path{RobroKools/vad-bert} descreve um modelo baseado em BERT ajustado para regressão das três dimensões e o identifica como recurso em inglês \cite{robrokoolsvadbert}. No VAD por LLM adotado neste estudo, as dimensões são definidas no \textit{prompt}, que solicita a avaliação do tom expresso em um texto individual, acompanhada de justificativas e trechos de apoio. As diferenças entre estimação lexical, neural e por LLM devem ser consideradas ao interpretar seus escores. Idioma, domínio, extensão processada e critérios de avaliação podem afetar a comparação, e a conversão para uma amplitude numérica comum não calibra os estimadores.

A dimensão afetiva também não esgota o enquadramento. A seleção e a saliência discutidas por \citeonline{entman1993framing} podem modificar responsabilização ou interpretação causal sem alterar substancialmente o tom. Duas mensagens com perfis VAD próximos podem apresentar posições opostas, assim como relatos de relações semelhantes podem expressar intensidades afetivas diferentes. Neste TCC, a análise afetiva complementa a leitura das transformações de conteúdo. O STDI não atribui um componente próprio à descrição textual de enquadramento, cuja função na representação é detalhada em Métodos.

As fontes examinadas fundamentam alternativas de representação, sem determinar quais estimadores são adequados a cada conjunto de textos. Os diagnósticos que orientaram a escolha do procedimento afetivo pertencem ao Capítulo \ref{ch:evolucao}, enquanto a configuração utilizada em cada corpus, a extração e as regras de agregação são descritas no Capítulo \ref{ch:metodos}.

\section{AVALIAÇÃO AUTOMÁTICA E POSICIONAMENTO DO ESTUDO}
\label{sec:avaliacao-llm}

A representação dos textos é uma etapa da mensuração. Outra é o procedimento que compara essas representações ou julga diretamente as versões. Um LLM pode atuar como gerador, extrator ou avaliador, funções que introduzem decisões em momentos diferentes. A expressão \textit{LLM-as-a-Judge} designa seu uso para produzir julgamentos segundo instruções. A nota obtida é uma estimativa condicionada ao material recebido e aos critérios definidos no \textit{prompt}.

Na comparação de versões, critérios separados permitem examinar o que motivou um julgamento. Uma avaliação de tema pode reconhecer continuidade do assunto, enquanto uma avaliação de relações pode identificar mudanças de causalidade, atribuição ou negação. A definição de níveis fornece âncoras para as notas e torna a decisão mais inspecionável. Essa organização não elimina a dependência do modelo e das instruções, nem estabelece que as notas correspondam a uma escala validada de mudança percebida por pessoas.

\citeonline{haldar2025ratingroulette} examinam a consistência de avaliadores LLM em diferentes tarefas de geração de linguagem e identificam variação entre execuções. O mesmo material pode, portanto, receber notas diferentes mesmo com modelo e instruções mantidos. \citeonline{saha2026judgebudget} estudam consultas repetidas para estimar a média das notas sob orçamento limitado e propõem sua alocação conforme a variância estimada dos pares. Esses trabalhos fundamentam a preocupação com repetibilidade, mas não estabelecem a estabilidade do STDI nem validam suas escolhas de agregação.

Em pesquisa social, \citeonline{abdurahman2025primer} enfatizam a documentação de modelos, versões e \textit{prompts}, a repetição de condições, a análise de sensibilidade e a comparação com julgamentos humanos. Essas recomendações ajudam a distinguir a possibilidade de recompor um procedimento da validade de sua interpretação. No presente estudo, a avaliação emprega comparações quantitativas, casos controlados e inspeções qualitativas. A ausência de referência humana numérica independente delimita as conclusões sobre os escores. Os procedimentos efetivamente realizados são descritos no Capítulo \ref{ch:metodos}.

A contribuição buscada pelo TCC está na articulação entre sequências controladas de perspectivas, representação estruturada e análise conjunta de mudanças locais e afastamento da origem. A decomposição do STDI permite examinar aspectos que um escore global isolado pode ocultar, enquanto a inspeção dos textos ajuda a interpretar divergências de mensuração. Essa proposta se insere na pesquisa sobre geração iterativa e agentes LLM, com alcance instrumental e descritivo. O capítulo seguinte apresenta as decisões que conduziram desse referencial à abordagem final.
